\documentclass[10pt,a4paper]{amsart}
\usepackage{amsmath,amssymb,amsthm,graphicx}
% Retained minimal preamble for the Stacks Project chapter more-morphisms.tex (More on Morphisms).
% Verbatim excerpt of the chapter preamble: Xy-pic setup (native preamble lines 14-18) and the native
% theorem declarations with the mathematical macros (native preamble lines 154-217). Copied mechanically;
% nothing added, removed or redefined. The wrapper class and the amsmath/amssymb/amsthm/graphicx base
% packages are supplied by the assembler, so the class/externaldocument part of the native preamble is omitted.
% For commutative diagrams we use Xy-pic
\usepackage[all]{xy}

% We use 2cell for 2-commutative diagrams.
\xyoption{2cell}
% Theorem environments.
%
\theoremstyle{plain}
\newtheorem{theorem}[subsection]{Theorem}
\newtheorem{proposition}[subsection]{Proposition}
\newtheorem{lemma}[subsection]{Lemma}

\theoremstyle{definition}
\newtheorem{definition}[subsection]{Definition}
\newtheorem{example}[subsection]{Example}
\newtheorem{exercise}[subsection]{Exercise}
\newtheorem{situation}[subsection]{Situation}

\theoremstyle{remark}
\newtheorem{remark}[subsection]{Remark}
\newtheorem{remarks}[subsection]{Remarks}

\numberwithin{equation}{subsection}

% Macros
%
\def\lim{\mathop{\mathrm{lim}}\nolimits}
\def\colim{\mathop{\mathrm{colim}}\nolimits}
\def\Spec{\mathop{\mathrm{Spec}}}
\def\Hom{\mathop{\mathrm{Hom}}\nolimits}
\def\Ext{\mathop{\mathrm{Ext}}\nolimits}
\def\SheafHom{\mathop{\mathcal{H}\!\mathit{om}}\nolimits}
\def\SheafExt{\mathop{\mathcal{E}\!\mathit{xt}}\nolimits}
\def\Sch{\mathit{Sch}}
\def\Mor{\mathop{\mathrm{Mor}}\nolimits}
\def\Ob{\mathop{\mathrm{Ob}}\nolimits}
\def\Sh{\mathop{\mathit{Sh}}\nolimits}
\def\NL{\mathop{N\!L}\nolimits}
\def\CH{\mathop{\mathrm{CH}}\nolimits}
\def\proetale{{pro\text{-}\acute{e}tale}}
\def\etale{{\acute{e}tale}}
\def\QCoh{\mathit{QCoh}}
\def\Ker{\mathop{\mathrm{Ker}}}
\def\Im{\mathop{\mathrm{Im}}}
\def\Coker{\mathop{\mathrm{Coker}}}
\def\Coim{\mathop{\mathrm{Coim}}}

% Boxtimes
%
\DeclareMathSymbol{\boxtimes}{\mathbin}{AMSa}{"02}

%
% Macros for moduli stacks/spaces
%
\def\QCohstack{\mathcal{QC}\!\mathit{oh}}
\def\Cohstack{\mathcal{C}\!\mathit{oh}}
\def\Spacesstack{\mathcal{S}\!\mathit{paces}}
\def\Quotfunctor{\mathrm{Quot}}
\def\Hilbfunctor{\mathrm{Hilb}}
\def\Curvesstack{\mathcal{C}\!\mathit{urves}}
\def\Polarizedstack{\mathcal{P}\!\mathit{olarized}}
\def\Complexesstack{\mathcal{C}\!\mathit{omplexes}}
% \Pic is the operator that assigns to X its picard group, usage \Pic(X)
% \Picardstack_{X/B} denotes the Picard stack of X over B
% \Picardfunctor_{X/B} denotes the Picard functor of X over B
\def\Pic{\mathop{\mathrm{Pic}}\nolimits}
\def\Picardstack{\mathcal{P}\!\mathit{ic}}
\def\Picardfunctor{\mathrm{Pic}}
\def\Deformationcategory{\mathcal{D}\!\mathit{ef}}
\newenvironment{reference}{}{}
\newenvironment{slogan}{}{}
\newenvironment{history}{}{}
\usepackage{cleveref}
\crefname{originallemma}{original lemma}{original lemmas}
\Crefname{originallemma}{Original lemma}{Original lemmas}
\makeatletter
\newlabel{algebra-lemma-flat-going-down}{{algebra-lemma-flat-going-down}{source}{}{sourcepointer-1}{}}
\newlabel{algebra-lemma-flat-going-down@cref}{{[originallemma][0][]algebra-lemma-flat-going-down}{[1][0][]source}{}{}{}}
\makeatother
\crefname{originallemma}{original lemma}{original lemmas}
\Crefname{originallemma}{Original lemma}{Original lemmas}
\makeatletter
\newlabel{algebra-lemma-integral-closure-commutes-colim-smooth}{{algebra-lemma-integral-closure-commutes-colim-smooth}{source}{}{sourcepointer-2}{}}
\newlabel{algebra-lemma-integral-closure-commutes-colim-smooth@cref}{{[originallemma][0][]algebra-lemma-integral-closure-commutes-colim-smooth}{[1][0][]source}{}{}{}}
\makeatother
\crefname{originallemma}{original lemma}{original lemmas}
\Crefname{originallemma}{Original lemma}{Original lemmas}
\makeatletter
\newlabel{algebra-lemma-integral-closure-commutes-etale}{{algebra-lemma-integral-closure-commutes-etale}{source}{}{sourcepointer-3}{}}
\newlabel{algebra-lemma-integral-closure-commutes-etale@cref}{{[originallemma][0][]algebra-lemma-integral-closure-commutes-etale}{[1][0][]source}{}{}{}}
\makeatother
\crefname{originallemma}{original lemma}{original lemmas}
\Crefname{originallemma}{Original lemma}{Original lemmas}
\makeatletter
\newlabel{algebra-lemma-minimal-prime-image-minimal-prime}{{algebra-lemma-minimal-prime-image-minimal-prime}{source}{}{sourcepointer-4}{}}
\newlabel{algebra-lemma-minimal-prime-image-minimal-prime@cref}{{[originallemma][0][]algebra-lemma-minimal-prime-image-minimal-prime}{[1][0][]source}{}{}{}}
\makeatother
\crefname{originallemma}{original lemma}{original lemmas}
\Crefname{originallemma}{Original lemma}{Original lemmas}
\makeatletter
\newlabel{more-algebra-lemma-fibres-henselization}{{more-algebra-lemma-fibres-henselization}{source}{}{sourcepointer-5}{}}
\newlabel{more-algebra-lemma-fibres-henselization@cref}{{[originallemma][0][]more-algebra-lemma-fibres-henselization}{[1][0][]source}{}{}{}}
\makeatother
\crefname{originallemma}{original lemma}{original lemmas}
\Crefname{originallemma}{Original lemma}{Original lemmas}
\makeatletter
\newlabel{more-algebra-lemma-henselization-noetherian}{{more-algebra-lemma-henselization-noetherian}{source}{}{sourcepointer-6}{}}
\newlabel{more-algebra-lemma-henselization-noetherian@cref}{{[originallemma][0][]more-algebra-lemma-henselization-noetherian}{[1][0][]source}{}{}{}}
\makeatother
\crefname{originallemma}{original lemma}{original lemmas}
\Crefname{originallemma}{Original lemma}{Original lemmas}
\makeatletter
\newlabel{morphisms-definition-normalization}{{morphisms-definition-normalization}{source}{}{sourcepointer-7}{}}
\newlabel{morphisms-definition-normalization@cref}{{[originallemma][0][]morphisms-definition-normalization}{[1][0][]source}{}{}{}}
\makeatother
\crefname{originallemma}{original lemma}{original lemmas}
\Crefname{originallemma}{Original lemma}{Original lemmas}
\makeatletter
\newlabel{morphisms-lemma-description-normalization}{{morphisms-lemma-description-normalization}{source}{}{sourcepointer-8}{}}
\newlabel{morphisms-lemma-description-normalization@cref}{{[originallemma][0][]morphisms-lemma-description-normalization}{[1][0][]source}{}{}{}}
\makeatother
\crefname{originallemma}{original lemma}{original lemmas}
\Crefname{originallemma}{Original lemma}{Original lemmas}
\makeatletter
\newlabel{morphisms-lemma-integral-closure}{{morphisms-lemma-integral-closure}{source}{}{sourcepointer-9}{}}
\newlabel{morphisms-lemma-integral-closure@cref}{{[originallemma][0][]morphisms-lemma-integral-closure}{[1][0][]source}{}{}{}}
\makeatother
\title{Normalization commutes with henselization and strict henselization at a point}
\author{The Stacks Project Authors}
\date{stacks-project:more-morphisms.tex:local-snapshot-20261011; GNU FDL 1.2 or later}
\begin{document}
\maketitle
\noindent\textit{Editorial note.} Statement span is lines 5438-5450 of sources/stacks-more-morphisms/pinned-original-source.tex; proof span is lines 5452-5496, the complete author proof (43 body lines).\par\medskip
\noindent\textit{Editorial note.} Navigation stubs: 9 cited-but-not-included labels are supplied as non-load-bearing pointers with their pinned tags where the snapshot tag index has them.\par\medskip
\noindent\textit{Editorial note.} extra\_packages is empty on purpose: the retained preamble already loads xy with options and a second declaration clashes in the assembler wrapper.\par\medskip
\noindent\textit{Editorial note.} Text is retained verbatim, including the author's own typos and placeholders.\par\medskip
\noindent\textit{Allowed standard prerequisites (editorial).}\begin{itemize}
\item Explicit prerequisite boundary. Normalization of a locally Noetherian scheme, integral closure, and the description of henselization and strict henselization as filtered colimits of etale algebras; the reductions are retained.
\end{itemize}\medskip
\ifdefined\hypertarget\hypertarget{sourcepointer-1}{}\fi
\noindent\textit{Source reference.} Original lemma algebra-lemma-flat-going-down: sources/stacks-more-morphisms/pinned-original-source.tex, source lines 5496--5496. Cited by the retained statement or proof but not part of the collected spans; supplied only as a navigation pointer. The argument does not depend on its content.\par\medskip
\ifdefined\hypertarget\hypertarget{sourcepointer-2}{}\fi
\noindent\textit{Source reference.} Original lemma algebra-lemma-integral-closure-commutes-colim-smooth: sources/stacks-more-morphisms/pinned-original-source.tex, source lines 5496--5496. Cited by the retained statement or proof but not part of the collected spans; supplied only as a navigation pointer. The argument does not depend on its content.\par\medskip
\ifdefined\hypertarget\hypertarget{sourcepointer-3}{}\fi
\noindent\textit{Source reference.} Original lemma algebra-lemma-integral-closure-commutes-etale: sources/stacks-more-morphisms/pinned-original-source.tex, source lines 5496--5496. Cited by the retained statement or proof but not part of the collected spans; supplied only as a navigation pointer. The argument does not depend on its content.\par\medskip
\ifdefined\hypertarget\hypertarget{sourcepointer-4}{}\fi
\noindent\textit{Source reference.} Original lemma algebra-lemma-minimal-prime-image-minimal-prime: sources/stacks-more-morphisms/pinned-original-source.tex, source lines 5496--5496. Cited by the retained statement or proof but not part of the collected spans; supplied only as a navigation pointer. The argument does not depend on its content.\par\medskip
\ifdefined\hypertarget\hypertarget{sourcepointer-5}{}\fi
\noindent\textit{Source reference.} Original lemma more-algebra-lemma-fibres-henselization: sources/stacks-more-morphisms/pinned-original-source.tex, source lines 5496--5496. Cited by the retained statement or proof but not part of the collected spans; supplied only as a navigation pointer. The argument does not depend on its content.\par\medskip
\ifdefined\hypertarget\hypertarget{sourcepointer-6}{}\fi
\noindent\textit{Source reference.} Original lemma more-algebra-lemma-henselization-noetherian: sources/stacks-more-morphisms/pinned-original-source.tex, source lines 5496--5496. Cited by the retained statement or proof but not part of the collected spans; supplied only as a navigation pointer. The argument does not depend on its content.\par\medskip
\ifdefined\hypertarget\hypertarget{sourcepointer-7}{}\fi
\noindent\textit{Source reference.} Original lemma morphisms-definition-normalization: sources/stacks-more-morphisms/pinned-original-source.tex, source lines 5496--5496. Cited by the retained statement or proof but not part of the collected spans; supplied only as a navigation pointer. The argument does not depend on its content.\par\medskip
\ifdefined\hypertarget\hypertarget{sourcepointer-8}{}\fi
\noindent\textit{Source reference.} Original lemma morphisms-lemma-description-normalization: sources/stacks-more-morphisms/pinned-original-source.tex, source lines 5496--5496. Cited by the retained statement or proof but not part of the collected spans; supplied only as a navigation pointer. The argument does not depend on its content.\par\medskip
\ifdefined\hypertarget\hypertarget{sourcepointer-9}{}\fi
\noindent\textit{Source reference.} Original lemma morphisms-lemma-integral-closure: sources/stacks-more-morphisms/pinned-original-source.tex, source lines 5496--5496. Cited by the retained statement or proof but not part of the collected spans; supplied only as a navigation pointer. The argument does not depend on its content.\par\medskip
\begin{lemma}[Normalization and henselization]
\label{lemma-normalization-and-henselization}
Let $X$ be a locally Noetherian scheme. Let $\nu : X^\nu \to X$
be the normalization morphism. Then for any point $x \in X$
the base change
$$
X^\nu \times_X \Spec(\mathcal{O}_{X, x}^h) \to \Spec(\mathcal{O}_{X, x}^h),
\quad\text{resp.}\quad
X^\nu \times_X \Spec(\mathcal{O}_{X, x}^{sh}) \to \Spec(\mathcal{O}_{X, x}^{sh})
$$
is the normalization of $\Spec(\mathcal{O}_{X, x}^h)$,
resp.\ $\Spec(\mathcal{O}_{X, x}^{sh})$.
\end{lemma}

\begin{proof}
Let $\eta_1, \ldots, \eta_r$ be the generic points of the irreducible
components of $X$ passing through $x$. The base change of the normalization to
$\Spec(\mathcal{O}_{X, x})$ is the spectrum of the
integral closure of $\mathcal{O}_{X, x}$ in $\prod \kappa(\eta_i)$.
This follows from our construction of the normalization of $X$ in
Morphisms, Definition \ref{morphisms-definition-normalization}
and Morphisms, Lemma \ref{morphisms-lemma-integral-closure};
you can also use the description of the normalization
in Morphisms, Lemma \ref{morphisms-lemma-description-normalization}.
Thus we reduce to the following algebra problem.
Let $A$ be a Noetherian local ring; recall that this implies
the henselization $A^h$ and strict henselization $A^{sh}$
are Noetherian too (More on Algebra, Lemma
\ref{more-algebra-lemma-henselization-noetherian}).
Let $\mathfrak p_1, \ldots, \mathfrak p_r$
be its minimal primes. Let $A'$ be the integral closure
of $A$ in $\prod \kappa(\mathfrak p_i)$.
Problem: show that
$A' \otimes_A A^h$, resp.\ $A' \otimes_A A^{sh}$ is constructed
from the Noetherian local ring $A^h$, resp.\ $A^{sh}$ in the same
manner.

\medskip\noindent
Since $A^h$, resp.\ $A^{sh}$ are colimits of \'etale $A$-algebras,
we see that the minimal primes of $A$ and $A^{sh}$ are exactly
the primes of $A^h$, resp.\ $A^{sh}$ lying over the minimal primes
of $A$ (by going down, see
Algebra, Lemmas \ref{algebra-lemma-flat-going-down} and
\ref{algebra-lemma-minimal-prime-image-minimal-prime}).
Thus More on Algebra, Lemma \ref{more-algebra-lemma-fibres-henselization}
tells us that
$A^h \otimes_A \prod \kappa(\mathfrak p_i)$,
resp.
$A^{sh} \otimes_A \prod \kappa(\mathfrak p_i)$
is the product of the residue fields at the minimal primes
of $A^h$, resp.\ $A^{sh}$. We know that taking the
integral closure in an overring commutes with \'etale
base change, see
Algebra, Lemma \ref{algebra-lemma-integral-closure-commutes-etale}.
Writing $A^h$ and $A^{sh}$ as a limit of \'etale $A$-algebras
we see that the same thing is true for the base change to
$A^h$ and $A^{sh}$ (you can also use the more general
Algebra, Lemma \ref{algebra-lemma-integral-closure-commutes-colim-smooth}).
\end{proof}

\end{document}
